% from ega2-1-fr.tex
\oldpage[II]{155}
une structure de $\mathcal{O}_X$-Algèbre graduée quasi-cohérente. En outre, le
diagramme
\[
\xymatrix{
  \mathcal{S}(h)\otimes_{\mathcal{S}}\mathcal{S}_+(k+1)\ar[r]\ar[d]_{1\otimes u_k}
  & \mathcal{S}_+(h+k+1)\ar[d]^{u_{k+h}}\\
  \mathcal{S}(h)\otimes_{\mathcal{S}}\mathcal{S}(k)\ar[r]
  & \mathcal{S}(h+k)
}
\]
est commutatif~; la fonctorialité des $\lambda$ montre alors que l'on a un
diagramme commutatif
\begin{equation}
\label{II.8.1.5.5-fr}
\xymatrix{
  \mathcal{O}_X(h)\otimes_{\mathcal{O}_X}\mathcal{O}_X(k+1)
    \ar[r]^{\lambda}\ar[d]_{1\otimes\widetilde{u}_k}
  & \mathcal{O}_X(h+k+1)\ar[d]^{\widetilde{u}_{k+h}}\\
  \mathcal{O}_X(h)\otimes_{\mathcal{O}_X}\mathcal{O}_X(k)
    \ar[r]^{\lambda}
  & \mathcal{O}_X(h+k)
}
\tag{8.1.5.5}
\end{equation}
les flèches horizontales étant les homomorphismes canoniques. On peut donc
dire que les $\widetilde{u}_k$ définissent un \emph{homomorphisme injectif} (de
degré $0$) \emph{de $\mathcal{S}_X$-Modules gradués}
\begin{equation}
\label{II.8.1.5.6-fr}
  \widetilde{u}:\mathcal{S}_X(1)\to\mathcal{S}_X.
\tag{8.1.5.6}
\end{equation}
\end{env}

\begin{env}[8.1.6]
\label{II.8.1.6-fr}
Gardant les notations de (\hyperref[II.8.1.5-fr]{8.1.5}), remarquons maintenant
que, pour $n\geq0$, l'homomorphisme composé
$\widetilde{v}_n=\widetilde{u}_{n-1}\widetilde{u}_{n-2}\cdots\widetilde{u}_0$
est un homomorphisme \emph{injectif} $\mathcal{O}_X(n)\to\mathcal{O}_X$~;
nous désignerons par $\mathcal{I}_{n,X}$ son image, qui est donc un Idéal
quasi-cohérent de $\mathcal{O}_X$, \emph{isomorphe} à $\mathcal{O}_X(n)$. En
outre, le diagramme
\[
\xymatrix{
  \mathcal{O}_X(m)\otimes_{\mathcal{O}_X}\mathcal{O}_X(n)
    \ar[r]^{\lambda}\ar[d]_{\widetilde{v}_m\otimes\widetilde{v}_n}
  & \mathcal{O}_X(m+n)\ar[d]^{\widetilde{v}_{m+n}}\\
  \mathcal{O}_X\ar[r]_{\mathrm{id}}
  & \mathcal{O}_X
}
\]
est commutatif pour $m\geq0$, $n\geq0$. On en conclut les inclusions suivantes
\begin{equation}
\label{II.8.1.6.1-fr}
  \mathcal{I}_{0,X}=\mathcal{O}_X,\qquad
  \mathcal{I}_{n,X}\subset\mathcal{I}_{m,X}
  \quad\text{pour }0\leq m\leq n
\tag{8.1.6.1}
\end{equation}
\begin{equation}
\label{II.8.1.6.2-fr}
  \mathcal{I}_{m,X}\mathcal{I}_{n,X}\subset\mathcal{I}_{m+n,X}
  \quad\text{pour }m\geq0,\ n\geq0.
\tag{8.1.6.2}
\end{equation}
\end{env}

